% ---------- Resumen ----------
\thispagestyle{plain}
\phantomsection\addcontentsline{toc}{chapter}{Resumen}
\begin{center}
  {\Large\bfseries RESUMEN}
\end{center}

\vspace{0.8\baselineskip}

\begin{spacing}{1.3}
    \justifying
    \setlength{\parskip}{6pt plus 1pt minus 1pt}
    \setlength{\parindent}{0pt}

La medición de peces a partir de imágenes evitaría retirarlos del agua, pero fotografiarlos a través de un tanque introduce un error: la luz se desvía al atravesar agua, vidrio y aire, y el tamaño aparente del pez depende de su distancia al vidrio, que una sola imagen no permite conocer. El objetivo fue comparar, sobre los mismos puntos anatómicos, cuatro procedimientos de conversión de píxeles a milímetros ---la escala de un marcador ArUco adherido al vidrio (A), una homografía con un marcador coplanar con la réplica (\Aprime), una calibración \emph{pinhole} bajo el agua (B) y un modelo refractivo de puerto plano (C)--- y cuantificar el intervalo de longitud que resulta de una profundidad desconocida. Se fotografiaron seis réplicas rígidas, medidas con calibre, a cuatro distancias del vidrio, con tres capturas por condición (72 imágenes). B y C requieren conocer la profundidad; con ese dato obtuvieron los menores errores de longitud total (MAE \SI{1,08}{mm}, \SI{0,73}{\percent}, y \SI{1,22}{mm}, \SI{0,88}{\percent}), seguidos de \Aprime{} (\SI{3,29}{mm}) y A (\SI{37,43}{mm}); no se detectó diferencia entre B y C (prueba de signos, $p = 1$). Sin información sobre la profundidad, el intervalo de C abarca del 56 al 82~\% de la longitud y contuvo la longitud de calibre en las 72 capturas; con la profundidad acotada a $\pm 1$~cm se reduce al 4--5~\%. Se concluye que, con la profundidad conocida, la calibración bajo el agua y el modelo refractivo reducen el error medio por debajo del 1~\%, y que sin ella la profundidad es la fuente dominante de incertidumbre.

    \textbf{Palabras clave:} \PalabrasClave.
\end{spacing}
\clearpage
